\section{Geometric Landauer: why correction costs are architecture-dependent}
\label{app:geometric_landauer}

Landauer's principle bounds the minimal dissipation for logically irreversible operations \cite{Landauer1961,Bennett1973}.
HPA--$\Omega$ refines this statement by emphasizing that operational information processing happens on constrained architectures (locality, routing, finite signal speed) and therefore incurs an additional geometric impedance cost \cite{Ma2025HPT,Ma2025CT}:
\begin{equation}
  W_{\mathrm{erase}}\ \ge\ \kB T_c\ln 2 + Z_{\mathrm{geom}}.
\end{equation}
The term $Z_{\mathrm{geom}}$ is not universal; it depends on the physical organization of the information-processing substrate.

\subsection{Biological proxies for geometric impedance}
\label{app:Zgeom_proxies}

In biological systems, candidate proxies for $Z_{\mathrm{geom}}$ include (non-exhaustively):
\begin{itemize}
  \item molecular transport distances and congestion (diffusion vs active transport),
  \item network topology and path-length distributions in signaling and regulatory networks,
  \item sparsity and long-range wiring costs in neural circuits,
  \item parallelism limits in repair pathways (bottlenecks and queueing),
  \item spatial localization constraints for assembly and proofreading operations.
\end{itemize}
Operationally, one can treat $Z_{\mathrm{geom}}$ as a fitted impedance term in an energy budget for correction tasks and then test whether its fitted variation correlates with measurable architectural features.


